\section{Fixed Points in Cartesian Closed Wild Categories}\label{sec:fixedpoints}

Taking a step back from the specific notions of this paper, we can consider abstractly the question of the existence of fixed points, and obtain a no-go theorem for the naïve setup where the kind ∗ is interpreted as \AgdaPrimitive{Type} and function types are preserved by the semantics, so that for instance the interpretation of a container $C$ of kind $(∗ → ∗) → ∗$ becomes an operation \AgdaSymbol{(}\AgdaPrimitive{Type}\;\AgdaSymbol{→}\;\AgdaPrimitive{Type}\AgdaSymbol{)}\;\AgdaSymbol{→}\;\AgdaPrimitive{Type}. Our chosen setting here is the notion of a cartesian closed wild category, and the naïve semantics would be a cartesian closed functor into the wild category of large types, sending the kind ∗ to \AgdaPrimitive{Type}.

\begin{definition}
A wild category $ℂ$, with $\Hom_ℂ(A,B)$ the type of morphisms from $A$ to $B$, is \emph{cartesian closed} if it is equipped with the following.

\begin{itemize}
\item
  A terminal object $1$: the type $\Hom_ℂ(C,1)$ is contractible for every object $C$. We write ${!} : C → 1$ for its unique morphism.
\item
  For all objects $A$ and $B$, a product $A×B$ with projections $π₁ : A×B → A$ and $π₂ : A×B → B$, such that $h ↦ (π₁ ∘ h, π₂ ∘ h)$ is an equivalence $\Hom_ℂ(C,A×B) ≃ \Hom_ℂ(C,A) × \Hom_ℂ(C,B)$ for every object $C$. We write $〈f,g〉 : C → A×B$ for the morphism corresponding to $f : C → A$ and $g : C → B$.
\item
  For all objects $A$ and $B$, an exponential $[A,B]$ with evaluation $\ev : [A,B]×A → B$, such that $h ↦ \ev ∘ 〈h ∘ π₁, π₂〉$ is an equivalence $\Hom_ℂ(C,[A,B]) ≃ \Hom_ℂ(C×A,B)$ for every object $C$. We write $Λ f : C → [A,B]$ for the morphism corresponding to $f : C×A → B$, and $g · a ≔ \ev ∘ 〈g,a〉 : C → B$ for the application of $g : C → [A,B]$ to $a : C → A$.
\end{itemize}
\end{definition}

A morphism $f : A → B$ has a \emph{name} $⌜f⌝ ≔ Λ(f ∘ π₂) : 1 → [A,B]$, where $π₂ : 1×A → A$. Conversely, a morphism $r : 1 → [A,B]$ determines the morphism $⌞r⌟ ≔ (r ∘ {!}) · \id : A → B$. These two operations are mutually inverse.

\begin{definition}
A \emph{functor} $F : ℂ → 𝔻$ of wild categories consists of an object $F␣A$ of $𝔻$ for every object $A$ of $ℂ$, and a morphism $F␣f : F␣A → F␣B$ for every morphism $f : A → B$, with identifications $F␣\id = \id$ and $F␣(g ∘ f) = F␣g ∘ F␣f$. A functor between cartesian closed wild categories is \emph{cartesian closed} if the following morphisms are isomorphisms, for all objects $A$ and $B$ of $ℂ$.

\begin{itemize}
\item
  The unique morphism ${!} : F␣1 → 1$.
\item
  The morphism $〈F␣π₁, F␣π₂〉 : F␣(A×B) → F␣A × F␣B$.
\item
  The morphism $Λ(F␣\ev ∘ 〈F␣π₁, F␣π₂〉^{-1}) : F␣[A,B] → [F␣A, F␣B]$.
\end{itemize}
\end{definition}

\begin{definition}
Let $f : A×B → B$ be a morphism in a cartesian closed wild category $ℂ$. An \emph{$A$-fixed point} of $f$ is a morphism $x : A → B$ together with an identification $f ∘ 〈\id, x〉 = x$.
\end{definition}

For an endomorphism $g : B → B$, a $1$-fixed point of $g ∘ π₂ : 1×B → B$ amounts to a morphism $x : 1 → B$ with an identification $g ∘ x = x$. In the wild category of large types, for instance, a $1$-fixed point of $F ∘ π₂$ for a function $F : \AgdaPrimitive{Type} → \AgdaPrimitive{Type}$ amounts by univalence to a type $X$ together with an equivalence $F␣X ≃ X$.

\begin{lemma}
\label{lem:fixed-preserved} A cartesian closed functor $F : ℂ → 𝔻$ maps $A$-fixed points in $ℂ$ to $F␣A$-fixed points in $𝔻$.
\end{lemma}

\begin{proof}
Let $x$ be an $A$-fixed point of $f : A×B → B$. We show that $F␣x$ is an $F␣A$-fixed point of $F␣f ∘ q$, where $q$ is any inverse of the canonical comparison map $F␣(A×B) → F␣A × F␣B$: \[F␣f ∘ q ∘ 〈\id, F␣x〉 = F␣f ∘ F␣〈\id, x〉 = F␣(f ∘ 〈\id, x〉) = F␣x.\]
\end{proof}

\begin{definition}
An \emph{internal fixed point operator} on an object $A$ of a cartesian closed wild category $ℂ$ is a morphism $y : [A,A] → A$ together with an identification $\id · y = y$.
\end{definition}

\[
\begin{tikzcd}[column sep=small]
  {[A,A]}
    \arrow[rr, "y", ""{name=Y, below}]
    \arrow[dr, "{〈\id, y〉}"'] &
  & A \\
  & {[A,A] × A}
    \arrow[ur, "\ev"']
    \arrow[to=Y, phantom, "="] &
\end{tikzcd}
\]

\begin{lemma}
Given an internal fixed point operator $y$ on $B$, every morphism $f : A×B → B$ yields an $A$-fixed point $y ∘ Λ f$.
\end{lemma}

\begin{proof}
\[f ∘ 〈\id, y ∘ Λ f〉 = \ev ∘ 〈Λ f, y ∘ Λ f〉 = \ev ∘ 〈\id, y〉 ∘ Λ f
  = y ∘ Λ f.\]
\end{proof}

Unfolding the notation, an internal fixed point operator on $A$ is exactly an $[A,A]$-fixed point of $\ev : [A,A]×A → A$.

\begin{proposition}
\label{prop:operator-preserved} A cartesian closed functor $F : ℂ → 𝔻$ maps internal fixed point operators on $A$ to internal fixed point operators on $F␣A$.
\end{proposition}

\begin{proof}
Let $y$ be an internal fixed point operator on $A$. We show that $F␣y ∘ m$ is an internal fixed point operator on $F␣A$, where $m$ is any inverse of the canonical comparison map $c : F␣[A,A] → [F␣A,F␣A]$: \[\ev ∘ 〈\id, F␣y ∘ m〉 = \ev ∘ 〈c, F␣y〉 ∘ m
= F␣\ev ∘ q ∘ 〈\id, F␣y〉 ∘ m = F␣y ∘ m.\] Here $q$ is as in the proof of Lemma \ref{lem:fixed-preserved}, so that $c = Λ(F␣\ev ∘ q)$, and the last step is that lemma applied to $y$ as an $[A,A]$-fixed point of $\ev$.
\end{proof}

\begin{definition}
For an object $A$ of a cartesian closed wild category, let $Φ_A : [[A,A],A] → [[A,A],A]$ be the morphism $Λ(π₂ · (π₁ · π₂))$ (which is $λ M f. f (M f)$, via the Curry--Howard--Lambek correspondence).
\end{definition}

\begin{lemma}
\label{lem:names} The map $⌜-⌝ : \Hom_ℂ([A,A],A) → \Hom_ℂ(1,[[A,A],A])$ is an equivalence, and lifts to an equivalence between the internal fixed point operators on $A$ and the $1$-fixed points of $Φ_A ∘ π₂$.
\end{lemma}

\begin{proof}
That $⌜-⌝$ is an equivalence was noted above. For the lift, it suffices that $Φ_A ∘ ⌜y⌝ = ⌜\id · y⌝$ for every $y : [A,A] → A$: \[Φ_A ∘ ⌜y⌝ = Λ(π₂ · ((⌜y⌝ ∘ π₁) · π₂)) = Λ(π₂ · (y ∘ π₂))
= Λ((\id · y) ∘ π₂) = ⌜\id · y⌝.\]
\end{proof}

\textbf{Remark:} Proposition \ref{prop:operator-preserved} can also be proved through $Φ_A$. Being given by a $λ$-term, $Φ_A$ is preserved by cartesian closed functors up to the comparison isomorphisms. So such a functor $F$ maps $1$-fixed points of $Φ_A ∘ π₂$ to $1$-fixed points of $Φ_{F␣A} ∘ π₂$, which by Lemma \ref{lem:names} are the internal fixed point operators on $F␣A$.

\begin{proposition}
A cartesian closed wild category with $1$-fixed points of $g ∘ π₂$ for all morphisms $g : [[A,A],A] → [[A,A],A]$ has an internal fixed point operator on $A$. I.e. an external fixed point operation of order $n+2$ gives rise to an internal fixed point operation of order $n$.
\end{proposition}

\begin{proof}
Take $g ≔ Φ_A$ and apply Lemma \ref{lem:names}.
\end{proof}

\begin{definition}
Let $\mathcal{T}$ be the wild category of large types: its objects are the types of the universe $\AgdaPrimitive{Type₁}$, among them $\AgdaPrimitive{Type}$, and its morphisms are functions. It is cartesian closed, with the unit type, product types and function types.
\end{definition}

\begin{lemma}
\label{lem:no-operator} $\mathcal{T}$ has no internal fixed point operator on $\AgdaPrimitive{Type}$.
\end{lemma}

\begin{proof}
Such an operator $y$ would give, for $D ≔ y␣(λ X. X → ⊥)$, an identification $(D → ⊥) = D$, from which $⊥$ is derivable.
\end{proof}

\begin{theorem}
Let $F : ℂ → \mathcal{T}$ be a cartesian closed functor and $∗$ an object of $ℂ$ with $F␣∗ ≃ \AgdaPrimitive{Type}$. Then $Φ_∗ ∘ π₂$ has no $1$-fixed point in $ℂ$.
\end{theorem}

\begin{proof}
Let $e : F␣∗ ≃ \AgdaPrimitive{Type}$, and suppose $r$ is a $1$-fixed point of $Φ_∗ ∘ π₂$. By Lemma \ref{lem:names} and Proposition \ref{prop:operator-preserved} there is an internal fixed point operator $y$ on $F␣∗$. Then $λ g. e␣(y␣(e^{-1} ∘ g ∘ e))$ is one on $\AgdaPrimitive{Type}$, contradicting Lemma \ref{lem:no-operator}.
\end{proof}



